\section{From acceptance probability to wire lower bounds}
\label{sec:lower}

Write $\THR[s]$ for real-weight, arbitrary-depth threshold circuits
with at most $s$ wires. Fixing or negating an input changes weights
and thresholds without increasing wires; output negation replaces
$[w\cdot y\ge\theta]$ by $[-w\cdot y\ge1-\theta]$ once integer weights
are chosen. Two-input AND, OR, and XOR have constant-wire implementations,
so an $r$-bit parity has $O(r)$ wires.
After a projection fixes every input of a gate, evaluate that gate and
replace it by constant unary tensors, preserving all remaining declared
inputs and adding no wires.

\subsection{Finite descriptions and the transfer ingredients}

A threshold function on $r$ bits has integer weights and threshold of
$O(r\log(r+2))$ bits. Here is the existence proof. Let $D=r+1$ and
for each $x$ put $a_x=(2f(x)-1)(x,-1)$. Move the original threshold
slightly downward to make all margins strict, and scale to obtain
$a_xz\ge1$ for all $x$. The rows have rank $D$, since the rows at
$0,e_1,\ldots,e_r$ do. A feasible point with maximal rank among its
tight rows has $D$ independent tight rows: otherwise move in a
nonzero direction orthogonal to them, choosing its sign to decrease
another inequality, until a new row becomes tight. This
increases the rank. For the resulting invertible matrix $B$, Cramer's
rule applied to $Bz=\mathbf1$ and Hadamard's inequality bound every
numerator and denominator by $D^{D/2}$. Multiplication by
$|\det B|$ gives the asserted integer representation. Constant
functions are handled separately.

Thus a $w$-wire threshold circuit has a normalized description of
$O((n+w)^2\log(n+w+2))$ bits. This remains true after projections.
A normalized fan-in-$r$ gate can also be simulated with
$O(r^2\log(r+2))$ fan-in-two Boolean gates: form the weighted sum
with binary addition and compare. A circuit with
$w=O(t\sqrt{\log t})$ and $n=O(t)$ consequently has ordinary
Boolean size at most $t^3$ for all sufficiently large $t$.

We use the following published ingredients in their specified forms.
They are quoted tools inside the proofs, and do not add an unproved
algorithmic assumption to \HTC.
All fixed probability thresholds used by a machine are chosen rational,
weakening completeness and soundness bounds within their positive gaps
when necessary. Thus the acceptance comparisons have finite encodings.
\begin{enumerate}
\item The unary nondeterministic time hierarchy supplies
$H\in\NTIME(2^t)\setminus\NTIME(2^{t/2})$
\cite[Theorem 2.1 (full version)]{Williams13}.
\item A succinct PCP for such computations has perfect completeness,
constant soundness $\beta<1$, proof and randomness indices padded to
$q=t+O(\log t)$, and polynomial-size query and verification circuits.
For the refuter computation of time $T(n)$ we use the stronger
soundness $n^{-10}$ and $q=\log_2T(n)+O(\log\log T(n))$ from
\cite[Lemma 3.10]{CLW20}. Its queries are projections. The projection
query lists and the verifier's fixed-width CNF give fixed polynomial
bounds in $q$ on the composed verifier size; their construction time
may additionally be polynomial in the original input length.
\item A constant-rate linear error-correcting code $\Enc$ with
polynomial-time decoder $\Dec$ has a fixed decoding radius. The
Circuit-Eval PCPP in \cite[Lemma 24]{CW19}, used in
\cite[Lemma 3.6]{BW24}, has perfect completeness, constant soundness
$\sigma<1$ at a fixed proximity, polynomial proof length and proof
construction, and constraints that are signed ORs of three queried
bits.
\item The refuter in \cite[Theorem 4.6]{CLW20} applies to any
time-constructible $n\le T(n)\le2^{\poly(n)}$. It supplies a language
$A^T_{\mathrm{FS}}\in\NTIME[T(n)]$ and, for each fixed
nondeterministic machine $M$ running in $o(T(n))$ time with at most
$n/10$ guessed bits, finds for every sufficiently large $n$ an input
$z\in\{0,1\}^n$ with $M(z)\ne A^T_{\mathrm{FS}}(z)$, in
$\poly(T(n))$ time with a SAT oracle.
\end{enumerate}

Fix once and for all the algorithm supplied by \HTC; it is part of
the finite description of each cheating machine below. The languages
constructed from their refuters are then fixed languages.
For each fixed description exponent $a$, fix an exponent constant
$A_5(a)$ such that the inverse-polynomial CAPP algorithm takes
\[
 \poly_a(S)2^{A_5(a)g(S)}
\]
time on normalized descriptions of length at most $S^a$, with error
at least $S^{-a}$. This constant is fixed before choosing wire-bound
constants.

All wire caps below can be made computable integer functions. For a
fixed positive rational $d$, use
\begin{equation}\label{eq:cap}
 B_d(r)=\left\lfloor 2d\,\operatorname{isqrt}
 \bigl(r^2\max\{1,\lfloor\log_2r\rfloor\}\bigr)\right\rfloor
 =(2d+o(1))r\sqrt{\log_2r}.
\end{equation}
The floor is computed using rational integer arithmetic. Every
candidate circuit is padded to a predetermined polynomial description
length, and invalid topologies or out-of-range codes are rejected.

\subsection{The infinitely-often construction}
\label{subsec:io}

Assume \HTC. Apply the succinct PCP to the unary hierarchy language
$H$. Fix a common nondecreasing index bound
$q(t)=t+O(\log t)$, and pad proof and randomness indices to this bound.
Extra random bits are ignored and unused proof bits are zero.
The lexicographically least proof accepted on all random strings,
if it exists, defines a proof-bit function $W_1(t,x)$; otherwise set
that function to zero. The full verifier has size $2^{O(t)}$ when all
random strings are written out. Prefix SAT queries recover its
lexicographically least proof in $2^{O(t)}$ time with an NP oracle.

Fix a small positive rational $d$. Define $W_2(t,x,j)$ as follows.
Find the lexicographically least normalized integer threshold circuit
$F_t$ of at most $B_d(q)$ wires whose truth table is a fully accepted
PCP proof. This search is in $2^{O(t)}$ time with SAT queries: for
each candidate description, write out the verification at every one
of the $2^q$ random strings, and regard their conjunction as an NP
condition on the description bits. Let $G_t(x)$ be rejection by the
original PCP using this circuit, and put
\[
 D_t(y)=1-G_t(\Dec(y)).
\]
The verifier $G_t$ is an ordinary Boolean circuit. If an
exact-threshold or MOD gate is used in such a verifier with
polynomial-bit parameters, expand it into binary arithmetic and
comparison circuits of polynomial size before applying the PCPP.
For every $x$ with $D_t(\Enc(x))=1$, take the canonical PCPP proof
from its polynomial-time constructor. Its $j$th bit is $W_2(t,x,j)$.
If no $F_t$ exists, or for an invalid index, output zero. Proof lengths
are padded to one fixed $t^u$ bound covering all legal candidates, so
$j$ has $\ell(t)=O(\log t)$ bits.

Both functions belong to $\ENP$ when inputs include a self-delimiting
$t$ header. Before doing any exponential work, check that $t$ is at
most the full input length and that the payload length is correct.
Form a single language $L_\star$ by a bit selecting $W_1$ or $W_2$.
Its legal lengths are $t+O(\log t)$, with possibly different padding
for the two branches. If $L_\star$ had circuits of at most
$dN\sqrt{\log_2N}$ wires at every sufficiently large $N$, projection
would give circuits for both bit functions within $B_d(q)$.

A fast nondeterministic simulator for $H(1^t)$ guesses $F$ and a
universal proof-bit circuit $T(x,j)$, each of at most $B_d(q)$ wires.
Generate the PCPP constraints $H_1,\ldots,H_m$ of the associated
$D$. Replace every code query by its parity on $x$ and every proof
query by the corresponding fixed-$j$ projection of the \emph{same}
$T$. The resulting threshold circuit $J_i(x)$ satisfies
\begin{equation}\label{eq:three-copy}
 w(J_i)\le3B_d(q)+O(q)
       \le6d\,q\sqrt{\log_2q}+O(q).
\end{equation}
For example, using six wires per two-input XOR gives the explicit
bound $3B_d(q)+18q+3$. The polynomial-size original verifier is used
to generate the PCPP; it is not copied into $J_i$.

If $H(1^t)=1$, the assumed small circuit for $W_1$ makes a legal
$F_t$ exist; the small circuit for $W_2$ supplies $T$, and every
$J_i$ is identically one. If $H(1^t)=0$, every $F$ makes $G(x)=1$
on at least $1-\beta$ of the inputs. For such $x$, every word within
the code's decoding radius of $\Enc(x)$ decodes to $x$, so is rejected
by $D$. The PCPP soundness then makes at least $1-\sigma$ of the
constraints fail for \emph{any} proof at that $x$. Averaging gives
some $i$ with
\[
 \Pr_x[J_i(x)=0]\ge\tau=(1-\beta)(1-\sigma)>0.
\]
CAPP error $1/S_i<\tau/4$ distinguishes these cases by testing all
estimates against $1-\tau/2$.

Integer normalization ensures $L_i\le S_i^5$ for all large $t$.
Set $A_5=A_5(5)$ before choosing $d$. From~\eqref{eq:three-copy},
\[
 S_i\le6d\,t\sqrt{\log_2t}+O(t),\qquad
 g(S_i)\le(6d+o(1))t.
\]
Choose $6A_5d<1/4$. Polynomially many deterministic queries and the
polynomial-size guesses then give a simulation below $2^{t/2}$ time,
contradicting the hierarchy. Thus infinitely many lengths require
more than $dN\sqrt{\log_2N}$ wires. This is the two-stage
Bathie--Williams proof instantiated at the present scale.

\subsection{Every sufficiently large length}
\label{subsec:ae}

Assume \HTC. Let $n$ denote the input length of the refuter, and
choose
\[
 T(n)=2^{\lceil n^{1/K}\rceil}
\]
for a fixed sufficiently large integer $K$. It is time constructible
and satisfies the refuter's range for all large $n$; finitely many
smaller values can be adjusted. At the sampled lengths $n=t^K$,
write $q(t)=t+O(\log t)$ for a fixed nondecreasing bound on PCP proof
indices and random bits. The two-stage simulator $A_{\mathrm{PCP}}$
just described now runs on arbitrary $z\in\{0,1\}^n$, using the PCP
for $A^T_{\mathrm{FS}}(z)$, with initial PCP soundness $n^{-10}$.
It guesses the two normalized threshold circuits $F,T$, each capped
by $B_d(q)$. The number of guessed bits is polynomial in $t$, with
fixed degree: the two circuit descriptions have that property even
though generating all PCPP constraints can take a larger polynomial
time. Choose $K$ greater than this degree. There are at most $n/10$
guessed bits for all large $n$.

The same error analysis shows
$A_{\mathrm{PCP}}(z)\le A^T_{\mathrm{FS}}(z)$. Its time is
\[
 \poly(n)+\poly(t)2^{(6A_5d+o(1))t}=o(2^t)
 \quad\text{if }6A_5d<1.
\]
The promise holds on every sufficiently large input length, with
$t=\lceil n^{1/K}\rceil$. The refuter therefore produces, for each
large sampled $n=t^K$, a $z_t$ satisfying
\begin{equation}\label{eq:refuter-gap}
 A^T_{\mathrm{FS}}(z_t)=1,
 \qquad A_{\mathrm{PCP}}(z_t)=0.
\end{equation}

It remains to encode both proof-bit functions at the \emph{same}
length. Let $h_0(t)=O(\log t)$ be a fixed canonical header length
and $\ell(t)=O(\log t)$ the uniformly padded proof-bit index length.
Use nondecreasing versions and define
\begin{equation}\label{eq:ae-padding}
 P(t)=h_0(t)+1+q(t)+\ell(t),\qquad
 t(N)=\max\{t:P(t)\le N\}.
\end{equation}
One may take $q(t)=t+C\lceil\log_2(t+2)\rceil$ and analogous
ceiling-log bounds for the other lengths. Thus $P(t+1)-P(t)=O(\log t)$
and $N=q(t)+O(\log t)$ for $t=t(N)$. Use fields
$(\text{header},b,x,j,\text{padding})$, with $|x|=q(t)$ and
$|j|=\ell(t)$. Invalid headers give zero. The $b=0$ branch outputs
the canonical fully accepted PCP proof bit $W_1(t,x)$ for $z_t$ and
ignores $j$; the $b=1$ branch uses the canonical budgeted $F_t$ and
outputs $W_2(t,x,j)$ as above. Both branches ignore padding.
The language is in $\ENP$: the refuter takes $2^{O(t)}$ time with
SAT queries, as do the explicitly expanded canonical searches, and
$t\le N$.

Suppose a circuit of at most $dN\sqrt{\log_2N}$ wires computes this
language at one sufficiently large $N$. Fix the correct header,
padding, and branch bit. For $W_1$ also fix $j$. The resulting
circuits have at most
\[
 dN\sqrt{\log_2N}\le B_d(q(t))
\]
wires for all large $N$, by the factor-two slack in~\eqref{eq:cap}.
The $W_1$ circuit makes a fully accepted $F_t$ exist; this ensures
that the second is a universal PCPP proof circuit for that canonical
$F_t$. Together they make $A_{\mathrm{PCP}}(z_t)=1$, contradicting
\eqref{eq:refuter-gap}. This works separately at every sufficiently
large $N$, proving the worst-case assertion of
Theorem~\ref{thm:lower}. Any sufficiently small positive rational
$d$ with $6A_5(5)d<1$ is admissible.

\subsection{Constant-advantage average-case hardness}
\label{subsec:average}

Assume \HTC. We give the polynomial-parameter version of the proof
in \cite[Section 6]{CLW20}, including the per-term size needed here.
Let the two-query PCPP in \cite[Lemma 3.11]{CLW20} have constants
$0<s_{\mathrm p}<c_{\mathrm p}<1$ and gap
$\Delta=c_{\mathrm p}-s_{\mathrm p}$. Its code bits are parities.
Its constraint count is a power of two, and its proof constructor is
polynomial time. Fix a rational
\begin{equation}\label{eq:delta0}
 0<\delta_0\le\Delta^2/10^6.
\end{equation}
Choose integer exponents $K_1,K_2,K_3$ as specified below, and then
choose the hierarchy's $K$ larger than the fixed degree of all guesses.

The cheating machine $A_{\mathrm{PCPP}}$ obtains the succinct PCP
for $A^T_{\mathrm{FS}}(z)$ and guesses an \emph{ordinary fan-in-two}
circuit $C_1$ of size $t^{K_1}$ as its oracle. The composed verifier
is an ordinary polynomial-size circuit. Apply the two-query PCPP,
with variables $Y\cup Z$, $m=\poly(t)$ constraints, and code word
$Y=\Enc(x)$. Enlarge fixed polynomial bounds on proof lengths and
constraint count uniformly across all candidate $C_1$; the native
power-of-two constraint convention makes binary indexing uniform.
To reach a common larger power of two, repeat every original
constraint equally many times, preserving all satisfaction averages.
The machine guesses, for each proof variable, a rational linear sum
of at most $M=t^{K_2}$ threshold circuits, each with at most $s_1$
wires, and total coefficient absolute value at most $t^{K_2}$.
Coefficients have a fixed $O(\log t)$-bit bound large enough for the
constructions below. These guesses are not assumed to take values
in $[0,1]$.

For a constraint $i$, let $T_{i1},T_{i2}$ be its two (unsigned)
variable functions. Let $F_i$ be its OR polynomial after replacing
a positive literal by $T$ and a negative one by $1-T$. Define
\begin{equation}\label{eq:validity-polys}
 P_{ij}=\begin{cases}
 T_{ij}^2(1-T_{ij})^2,&T_{ij}\text{ corresponds to a }Z\text{ variable},\\
 (\Enc_s-T_{ij})^2,&T_{ij}\text{ corresponds to }Y_s,
 \end{cases}
 \qquad Q_{ij}=T_{ij}^2.
\end{equation}
Estimate expectations of all these polynomials by expanding products
of at most four sum terms. Each expectation is a linear combination
of at most polynomially many expectations of ANDs of at most four
threshold circuits, together with an $O(q)$-wire parity when required.
Every CAPP query has
\begin{equation}\label{eq:four-copy}
 w\le4s_1+O(q).
\end{equation}
The parity is computed as part of the circuit; no enumeration of its
input support is needed. Taking per-product error $t^{-K_3}$ gives
total expectation error $O(t^{4K_2-K_3})=o(1)$ when $K_3>4K_2$.
Increase $K_3$ so all errors below are eventually at most $\delta_0/10$.

Reject unless the estimated average of $P_{ij}$ is at most
$2\delta_0$ and every estimated $Q_{ij}$ is at most $1+\delta_0$.
After this validity test, accept precisely when
\begin{equation}\label{eq:avg-accept}
 \frac1m\sum_i\widetilde{\E}_x F_i
 \ge c_{\mathrm p}-\Delta/2.
\end{equation}
These are the tests of \cite[Section 6.2.2]{CLW20}; their requisite
estimates involve only~\eqref{eq:four-copy}.

Here are the validity estimates explicitly. If the guessed functions
are $[0,1]$-valued and their average $\ell_1$ distance, over occurrences,
from a Boolean PCPP proof is at most $\delta_0$, then the true average
of $P$ is at most $\delta_0$ and every $Q$ is at most one. Thus they
pass. Conversely, passing implies a true average $P\le3\delta_0$
and $\E Q_{ij}\le1+2\delta_0$. Round every $Z$ function to its
nearest Boolean value, and set $Y=\Enc(x)$. For real $u$,
\[
 \min\{u^2,(1-u)^2\}\le4u^2(1-u)^2.
\]
Hence, by Cauchy--Schwarz, the average $L_2$ distance of the guessed
occurrences from the rounded proof is at most
$\sqrt{12\delta_0}$. For each OR polynomial, telescoping its two
arguments, using the Boolean rounded argument and
$\|T_{ij}\|_2\le\sqrt{1+2\delta_0}$, gives
\[
 \|F_i-\widehat F_i\|_1
 \le7\,\E_{j\in\{1,2\}}\|T_{ij}-\widehat T_{ij}\|_2.
\]
For a negated literal use
$\|1-T_{ij}\|_2\le1+\sqrt{1+2\delta_0}$ in the same telescoping;
this gives the constant $7$. If
$A^T_{\mathrm{FS}}(z)=0$, the initial PCP accepts on at most $n^{-10}$
of the $x$'s, so every Boolean PCPP proof has average satisfaction
at most $s_{\mathrm p}+o(1)$. The guessed proof's average is at most
\[
 s_{\mathrm p}+7\sqrt{12\delta_0}+o(1)
 <c_{\mathrm p}-\Delta/2.
\]
Including estimation error preserves the strict inequality. Therefore
$A_{\mathrm{PCPP}}(z)\le A^T_{\mathrm{FS}}(z)$ for all large inputs.

Now fix $\gamma\in(0,1/2)$, and choose the least $k\ge1$ with
\begin{equation}\label{eq:xor-k}
 \eps_k=(1-\delta_0)^{k-1}(1/2-\delta_0)\le\gamma/2.
\end{equation}
Let $c_{\mathrm{STV}}$ be a fixed integer upper bound on the input-length blow-up in the
ordinary-circuit amplification \cite[Lemma 3.9]{CLW20}, originating
in~\cite{STV01}, and put
\begin{equation}\label{eq:kprime}
 k'=\max\{k,c_{\mathrm{STV}}\}.
\end{equation}
Take $s_1=B_{dk}(q)$. Fix
$a_*\ge\max\{5,K_3\}$ large enough for all query descriptions and
precisions, and write $A_5=A_5(a_*)$. Choose a small positive rational
$d$ such that
\begin{equation}\label{eq:avg-time-condition}
 8A_5dk<1.
\end{equation}
Indeed,~\eqref{eq:four-copy} has $4s_1\le8dkq\sqrt{\log_2q}$, so
the machine's time is
$\poly(t)2^{(8A_5dk+o(1))t}=o(2^t)$. The fixed exponent $a_*$ is
chosen to cover the required inverse-polynomial error; the constant
$A_5$ need not equal $A_5(5)$ from the worst-case argument.
All guesses have polynomial length in $t$, so a sufficiently large
$K$ gives the refuter promise. Obtain $z_t$ with
$A^T_{\mathrm{FS}}(z_t)=1$ and $A_{\mathrm{PCPP}}(z_t)=0$.

There are two cases, decidable in $2^{O(t)}$ time with SAT queries by
expanding all PCP random strings.
If no $C_1$ of size $t^{K_1}$ gives a fully accepted oracle, the
canonical fully accepted proof's truth table has ordinary circuit
complexity greater than $t^{K_1}$. The amplification in
\cite[Lemma 3.9]{CLW20} constructs, in $2^{O(q)}$ time from that
truth table, a function $g$ on at most $c_{\mathrm{STV}}q$ inputs
that is $(1/2+t^{-1})$-hard for ordinary size $t^3$, provided
$K_1$ exceeds three times the fixed amplification exponent. The
Boolean simulation of threshold circuits above makes $g$ hard for
the required near-linear threshold wire budget as well.

In the other case take the lexicographically least successful ordinary $C_1$
and a canonical Boolean PCPP proof $\widehat T$, using the proof
constructor separately for each $x$. Set
\[
 f(i,j,x)=\widehat T_{ij}(x),\qquad
 n'=q+\log_2m+1=t+O(\log t).
\]
This function is $\delta_0$-far in $\ell_1$ from every $[0,1]$-valued
linear sum with at most $M/(2m)$ terms, each of at most $s_1$ wires,
and coefficient absolute sum at most $t^{K_2}/(2m)$. To see this,
suppose $D$ is a closer sum and define the guessed function for each
proof variable $X$ by averaging $D(i,j,x)$ over occurrences of $X$.
Empty occurrence lists can be assigned zero. Each resulting function
has at most $M$ terms, no larger per-term wire bound, and no larger
coefficient absolute sum than the allowed cap. For $v_i\in[0,1]$
and $b\in\{0,1\}$,
\[
 \left|\frac1h\sum_i v_i-b\right|=
 \frac1h\sum_i|v_i-b|.
\]
Apply this equality for each $x$ and occurrence list: the average
$\ell_1$ distance to $\widehat T$ remains at most $\delta_0$.
The validity test passes, and the OR polynomials, Lipschitz on
$[0,1]^2$, lose at most $2\delta_0$ average satisfaction relative
to the canonical proof. They therefore satisfy~\eqref{eq:avg-accept},
contradicting $A_{\mathrm{PCPP}}(z_t)=0$.

Apply the \emph{proof of Lemma 3.8 in Appendix A} of~\cite{CLW20},
rather than only its complexity statement. For completeness, its
reconstruction has the following form. Given a circuit approximating
$f^{\oplus k}$, either fix one block and recurse on $k-1$, possibly
negating the circuit, or sample
$\ell=O(n'/(\delta_0\eps_k)^2)$ fixed values of the other blocks.
A Chernoff bound and a union bound over the $2^{n'}$ remaining inputs
make the empirical conditional agreement uniformly accurate. An
affine rescaling of that empirical mean produces a $[0,1]$-valued
linear sum approximating $f$ within $\delta_0$ in $\ell_1$.
Every term is a projection, possibly negated, of the original
approximator. In particular, each term retains its original wire
bound; there are $O_\gamma(t)$ terms, and the sum of coefficient
absolute values is $O_\gamma(1)$, or the weaker sufficient
$O_\gamma(t)$. All coefficients are rational with $O_\gamma(\log t)$
bits because $\delta_0$ and $\eps_k$ are rational. Increasing the
fixed $K_2$ makes these sums fit the preceding caps. Thus
$f^{\oplus k}$ cannot have agreement $1/2+\eps_k$ with the
threshold approximators under consideration.

Finally use a nondecreasing payload bound
\[
 P_\gamma(t)=h_0(t)+
 \max\{k(q(t)+\log_2m(t)+1),\ c_{\mathrm{STV}}q(t)\}
 =(k'+o(1))t,
\]
where polynomial padding of the PCPP supplies a common dyadic
$m(t)$ bound. For each input length $N$ choose the largest $t$ with
$P_\gamma(t)\le N$, and read the active function input from the leading payload bits, ignoring
all other bits. The length $N$ determines $t$; no header is validated and
every word of length $N$ is used. In Case 1, output $g$; in Case 2, output $f^{\oplus k}$. Canonical
searches, case selection, and truth-table amplification all run in
$2^{O(t)}$ time with an NP oracle, so this defines $L_\gamma\in\ENP$.
Set $c_\gamma=dk/k'$. A length-$N$ approximator of at most
$c_\gamma N\sqrt{\log_2N}$ wires has
\[
 s_0\le(dk+o(1))t\sqrt{\log_2t}<B_{dk}(q)=s_1
\]
wires for all large $t$. Averaging over padding gives a fixing with
the same agreement. Both cases yield a contradiction to agreement
greater than $1/2+\gamma$. One can take
\[
 c_\gamma=\Theta\!\left(\frac{1}{A_5k'}\right)
 =\Theta\!\left(\frac{\delta_0}{A_5\log(1/\gamma)}\right)
 \quad\text{as }\gamma\downarrow0.
\]
This proves the average-case assertion of
Theorem~\ref{thm:lower}. The proof uses inverse-polynomial CAPP
precision throughout.
